% ---------------------------------------------------------------
%  Muc II.2 -- Thuc trang (tinh trang giai phap da biet)
% ---------------------------------------------------------------
\subsection{Thực trạng vấn đề --- tình trạng các giải pháp đã biết}

\subsubsection{Thực trạng biên soạn đề kiểm tra tại đơn vị}

Để có căn cứ khách quan, tôi đã khảo sát giáo viên tổ Toán của nhà trường
trước khi triển khai sáng kiến.

\begin{center}
\begin{tabular}{|c|p{8.2cm}|c|c|}
\hline
\textbf{TT} & \textbf{Nội dung khảo sát} & \textbf{Số GV} & \textbf{Tỉ lệ} \\ \hline
1 & Thời gian soạn một đề kiểm tra định kì trên 4 giờ & & \\ \hline
2 & Cho rằng khó bảo đảm đúng tỉ lệ bốn mức độ nhận thức & & \\ \hline
3 & Tạo mã đề bằng cách đảo thứ tự câu hỏi và phương án & & \\ \hline
4 & Không tái sử dụng được kho câu hỏi của năm trước & & \\ \hline
5 & Chưa quen soạn đủ bốn dạng câu hỏi theo cấu trúc mới & & \\ \hline
\end{tabular}
\end{center}
\textit{Bảng 1. Kết quả khảo sát giáo viên tổ Toán (tổng số: \dots\ giáo viên)}

\ghichu{Chị phát phiếu cho tổ Toán rồi điền hai cột cuối. Nếu tổ ít người, ghi
thẳng số lượng (ví dụ 5/7) --- hội đồng không đòi mẫu lớn, chỉ đòi có số thật.
Nếu chưa kịp khảo sát, có thể bỏ bảng này và viết thành đoạn văn nhận định,
nhưng có bảng thì thuyết phục hơn nhiều.}

Qua trao đổi trực tiếp, các giáo viên đều cho rằng khâu tốn công nhất không
phải là nghĩ ra câu hỏi mà là \textit{cân đối ma trận} và \textit{tạo các mã đề
khác nhau}. Một đề kiểm tra định kì theo cấu trúc mới gồm 12 câu trắc nghiệm
nhiều lựa chọn, 2 câu đúng/sai (tương đương 8 ý), 4 câu trả lời ngắn và 3 câu
tự luận; để có bốn mã đề thực sự khác nhau thì phải soạn tới bốn đề độc lập ---
một khối lượng công việc không khả thi trong điều kiện dạy học bình thường.

\subsubsection{Thực trạng tự luyện tập của học sinh}

\begin{center}
\begin{tabular}{|c|p{8.2cm}|c|c|}
\hline
\textbf{TT} & \textbf{Nội dung khảo sát} & \textbf{Số HS} & \textbf{Tỉ lệ} \\ \hline
1 & Có nhu cầu tự luyện thêm bài tập sau mỗi bài học & & \\ \hline
2 & Tìm đề luyện tập trên mạng, không rõ có đúng chương trình không & & \\ \hline
3 & Gặp lời giải sai hoặc không hiểu được khi tự tra trên mạng & & \\ \hline
4 & Muốn biết ngay đúng/sai sau khi làm bài & & \\ \hline
5 & Không biết mình đang yếu ở bài nào cụ thể & & \\ \hline
\end{tabular}
\end{center}
\textit{Bảng 2. Kết quả khảo sát học sinh trước khi áp dụng (tổng số: \dots\ học sinh)}

\ghichu{Khảo sát nhanh bằng Google Biểu mẫu trong 5 phút đầu tiết là đủ. Nhớ
giữ lại ảnh chụp biểu mẫu và bảng kết quả để đưa vào phụ lục --- đó là bằng
chứng tốt.}

Điểm đáng chú ý nhất là học sinh không thiếu \textit{ý chí} luyện tập mà thiếu
\textit{nguồn đề đúng tầm}: đề trên mạng hoặc quá khó, hoặc thuộc chương trình
cũ, hoặc có lời giải sai; và dù có làm thì cũng phải chờ giáo viên chữa mới
biết mình đúng hay sai.

\subsubsection{Phân tích các giải pháp đã có}

Trên thực tế hiện đã có ba nhóm giải pháp, mỗi nhóm giải quyết được một phần và
để lại một hạn chế cốt lõi.

\textbf{Nhóm 1 --- Ngân hàng đề tĩnh do giáo viên tự tập hợp.} Giáo viên sưu
tầm câu hỏi từ nhiều nguồn, gõ vào Word hoặc \LaTeX\ rồi cắt ghép thành đề.
Ưu điểm là nội dung do chính giáo viên kiểm soát. Hạn chế cốt lõi: kho câu hỏi
là \textit{hữu hạn và tĩnh} --- dùng một lần là học sinh biết đề, muốn có đề
mới phải ngồi thay số bằng tay từng câu một, và mỗi lần thay số lại phải giải
lại để kiểm tra đáp án.

\textbf{Nhóm 2 --- Phần mềm trộn đề và các website ngân hàng đề thương mại.}
Các phần mềm này lấy câu hỏi từ một kho có sẵn của nhà cung cấp rồi hoán vị thứ
tự câu và thứ tự phương án. Hạn chế cốt lõi: các mã đề chỉ khác nhau về
\textit{thứ tự}, nội dung vẫn y hệt; giáo viên thường không đưa được ngân hàng
câu hỏi của riêng mình vào; và kho câu hỏi vẫn là kho tĩnh, chỉ là tĩnh của
người khác.

\textbf{Nhóm 3 --- Dùng trí tuệ nhân tạo sinh thẳng đề.} Đây là hướng đang được
nói đến nhiều nhất hiện nay: đưa yêu cầu cho một mô hình ngôn ngữ lớn và nhận
về một đề hoàn chỉnh. Ưu điểm là nhanh và không cần chuẩn bị gì. Hạn chế cốt
lõi nằm ở chính bản chất của mô hình --- nó dự đoán chuỗi kí tự có xác suất cao
nhất chứ không thực hiện phép toán, nên trong quá trình thử nghiệm tôi thường
xuyên gặp các lỗi sau:

\begin{itemize}[leftmargin=1.2cm, itemsep=2pt]
  \item \textbf{Vượt khung chương trình:} đưa vào kiến thức của lớp trên hoặc
        của chương chưa học, dù yêu cầu đã ghi rõ phạm vi.
  \item \textbf{Tính toán sai:} đáp án không khớp với đề bài, hoặc lời giải
        trình bày trôi chảy nhưng sai ở một bước biến đổi.
  \item \textbf{Không ổn định:} cùng một yêu cầu, chạy hai lần cho hai kết quả
        khác nhau, không kiểm chứng lại được.
  \item \textbf{Sai mức độ:} câu ghi là nhận biết nhưng thực chất là vận dụng,
        làm hỏng ma trận đề.
\end{itemize}

Một đề kiểm tra lấy điểm chính thức không chấp nhận được các rủi ro đó.

\begin{center}
\begin{tabular}{|p{3.2cm}|c|c|c|c|}
\hline
\textbf{Tiêu chí} & \textbf{Ngân hàng tĩnh} & \textbf{Phần mềm trộn đề} &
\textbf{AI sinh đề} & \textbf{Sáng kiến này} \\ \hline
Nội dung do giáo viên quyết định & Có & Không & Không & Có \\ \hline
Sinh được đề mới, khác số liệu & Không & Không & Có & Có \\ \hline
Các mã đề khác nhau thật sự & Không & Không & Có & Có \\ \hline
Đáp án, lời giải chắc chắn đúng & Có & Có & Không & Có \\ \hline
Đúng khung chương trình & Có & Tuỳ nguồn & Không chắc & Có \\ \hline
Đúng ma trận bốn mức độ & Thủ công & Tuỳ nguồn & Không chắc & Tự động \\ \hline
\end{tabular}
\end{center}
\textit{Bảng 3. So sánh sáng kiến với các giải pháp đã có}

Bảng trên cho thấy khoảng trống mà sáng kiến này lấp vào: \textbf{vừa giữ được
quyền quyết định nội dung của giáo viên như ngân hàng tĩnh, vừa sinh được đề
mới như trí tuệ nhân tạo, mà không mang theo rủi ro sai sót của trí tuệ nhân
tạo}. Cách làm để đạt được cả ba điều đó cùng lúc là: câu hỏi do chính giáo
viên viết ra, nhưng viết dưới dạng \textit{chương trình Python sinh câu hỏi}
chứ không phải dưới dạng văn bản cố định. Toàn bộ số liệu, đáp án và lời giải
đều do Python tính ra bằng công thức toán học, nên không thể sai; phạm vi kiến
thức đã bị khoá cứng ngay trong mã định danh của câu hỏi, nên không thể vượt
khung chương trình.

\clearpage
